\section{Related Work}
\label{sec:related_work}
The application of artificial intelligence to climate science has evolved rapidly. Flood prediction frameworks have increasingly relied upon ensemble tree architectures, such as Random Forests, owing to their non-linear discrimination capabilities and interpretability across diverse hydrological catchments \cite{Mosavi2018}. Concurrently, gradient boosting architectures like LightGBM \cite{Ke2017} and XGBoost \cite{Chen2016} have established dominance in evaluating localized temperature extremes and heatwave events, frequently outperforming traditional thermodynamic numerical models over short-to-medium range horizons.

Beyond isolated events, the prediction of compound climate hazards---where multiple extremes co-occur or sequentially interact---has emerged as a critical research frontier \cite{Zscheischler2018}. Deep learning techniques, particularly Recurrent Neural Networks (RNNs) utilizing LSTM and GRU cells, are highly proficient at mapping these complex temporal interactions \cite{Chattopadhyay2020}, \cite{Cho2014}. However, literature assessing the operationalization of these architectures frequently overlooks the stringent computational requirements imposed by maintaining live, uninterrupted historical sequence buffers during API-driven realtime inference. Our framework addresses this gap by directly comparing the theoretical offline capacity of sequence models against the practical viability of a stateless stacking ensemble for compound risk generation.
